\section{Intersection cohomology of the theta divisor}\label{sec:intersection}

Let \(j:U\hookrightarrow\Theta\).  For traditional middle perversity,
\(\bar m(8)=3\), and the unshifted Deligne construction for this isolated
singularity is
\[
 \mathcal P_{\bar m}=\tau_{\leq3}Rj_*\Z_U;
\]
see \cite[Section~3.1 and Theorem~3.5]{GoreskyMacPhersonIHII}.  Hence its
degree-three hypercohomology is ordinary restriction to the smooth locus,
so
\begin{equation}\label{eq:ih-smooth-locus}
 IH^3(\Theta,\Z)\cong H^3(U,\Z).
\end{equation}
The \(\Z/3\) term in the link occurs in degree four by
Lemma~\ref{lem:link-cohomology}; it lies beyond this truncation and does not
enter \eqref{eq:ih-smooth-locus}.
The Mayer--Vietoris proof of Proposition~\ref{prop:topological-exact}
identifies \(H^3(M,\Z)\to H^3(U,\Z)\) with an isomorphism.  This proves the
first assertion of Corollary~\ref{cor:main-ih}.  Under that identification,
the pair sequence \eqref{eq:theta-pair} and the surjectivity supplied by
Proposition~\ref{prop:endpoint-lifts} give
\[
 0\longrightarrow H^3(\Theta,\Z)
 \longrightarrow IH^3(\Theta,\Z)
 \longrightarrow H^3(X,\Z)\longrightarrow0,
\]
which proves the corollary.

For comparison, the rational decomposition of this resolution is already
known.  Krämer proves, with complex coefficients, the following noncanonical
derived-category decomposition \cite[Corollary~6]{KraemerE6}:
\begin{equation}\label{eq:kraemer-decomposition}
 R\sigma_*\Q_M[4]
 \cong\IC_\Theta\oplus\Q_0[2]\oplus\Q_0\oplus\Q_0[-2],
\end{equation}
where the same argument works over \(\Q\).  Equivalently,
\[
 H^k(M,\Q)\cong IH^k(\Theta,\Q)\oplus
 \begin{cases}
  \Q,&k=2,4,6,\\
  0,&\text{otherwise}.
 \end{cases}
\]
The three skyscraper summands come from the algebraic cohomology of the
exceptional cubic.  Equation \eqref{eq:kraemer-decomposition} does not split
the integral sequence in a manner compatible with \((b_*,e^*)\), nor does
it determine the mod-two coset in \eqref{eq:main-fibre-product}.

The integral calculation outside degree three has an additional coefficient
issue: in the link Gysin sequence, multiplication by the hyperplane class
sends \(h\) to \(3\ell\).  The next section identifies this map as the
central integral obstruction hidden by \eqref{eq:kraemer-decomposition}.
